\section{Sequence Diagrams}

Each sequence diagram has the user first, then the boundary, control, data and
entity classes from left to right. Messages are numbered and use the operation
names of the class diagram. Activation bars show when an object is active and
dashed arrows are returns. The fragments \texttt{alt} and \texttt{loop} show
alternatives and repetition.

\subsection*{Find available flights (UC-04), control class FlightSearchController}

The customer enters the criteria in FlightSearchView, which calls
\texttt{search(criteria)} on FlightSearchController. The controller validates
the criteria itself. If they are invalid it returns an error and the view shows
it. Otherwise it asks FlightDao for the flights of that day and, for each
flight, asks Flight whether it has departed and SeatDao how many seats are
free. The controller returns the offers, and the view shows a message when
there are none and the table of flights when there are some.

\diagramland{30-sequence-find-flights}{Sequence diagram: find available flights (UC-04)}{fig:seqfind}

\subsection*{Book a flight (UC-06), control class BookingController}

The customer chooses Book Flight and SeatSelectionView calls
\texttt{startReservation(customerId, flightId)} on BookingController. The
controller creates a Reservation, whose state is \texttt{CREATED}, and saves it
through ReservationDao. The view then asks for \texttt{availableSeats(flightId)}
and shows the seat map. When the customer picks a seat, the view calls
\texttt{selectSeat(reservationId, seatNo)}. The controller loads the
reservation, calls \texttt{assignSeat(seatNo, price)}, which moves it to
\texttt{PENDING\_PAYMENT}, and updates it. If another customer took the seat
first, ReservationDao raises \texttt{SeatUnavailableException} and the view
shows a message and a fresh seat map. Otherwise the view shows the price and
the Make Payment process takes over, which is shown by the reference frame.

\diagramland{31-sequence-book-flight}{Sequence diagram: book a flight (UC-06)}{fig:seqbook}

% The sequence diagrams for UC-07 and UC-14 are added here by issues P3-07 and P3-08.
